\documentclass[a4paper,landscape,headrule,footrule]{foils}


%%
%%% macros for Theories of Grammar
%%%
\usepackage{polyglossia}
\setdefaultlanguage{english}
%\setmainfont{TeX Gyre Pagella}


\newcommand{\logo}{~}
\newcommand{\header}[3]{%
\title{\vspace*{-2ex} \large JMORF/MORS --- Morpho-Syntax
\\[2ex] \Large  \emp{#2} \\ \emp{#3}}
\author{\blu{Francis Bond}   \\ 
\normalsize  \textbf{Palacký University}\\
\normalsize  \url{https://fcbond.github.io/}\\
\normalsize  \texttt{bond@ieee.org}}
\MyLogo{JMORF (2026)}
\renewcommand{\logo}{#2}
\hypersetup{
   pdfinfo={
     Author={Francis Bond},
     Title={#1: #2},
     Subject={JMORF/MORS: Morpho-Syntax},
     Keywords={Syntax, Morphology, HPSG, Unification, Constructions},
     License={CC BY 4.0}
   }
 }

\date{#1
  \\ Location: KB-1.11}
}
\usepackage{hyperref}
\hypersetup{
     colorlinks,
     linkcolor={blue!50!black},
     citecolor={red!50!black},
     urlcolor={blue!80!black}
   }



\usepackage{xcolor}
\usepackage{graphicx}
\newcommand{\blu}[1]{\textcolor{blue}{#1}}
\newcommand{\grn}[1]{\textcolor{green}{#1}}
\newcommand{\hide}[1]{\textcolor{white}{#1}}
\newcommand{\emp}[1]{\textcolor{red}{#1}}
\newcommand{\txx}[1]{\textbf{\textcolor{blue}{#1}}}
\newcommand{\lex}[1]{\textbf{\mtcitestyle{#1}}}


% \renewcommand{\labelitemi}{\textcolor{violet}{‣}}
% \renewcommand{\labelitemii}{\textcolor{purple}{‣}}

\newcommand{\subhead}[1]{\noindent\textbf{#1}\\[5mm]}

\newcommand{\Bad}{\emp{\raisebox{0.15ex}{\ensuremath{\mathbf{\otimes}}}}}
\newcommand{\bad}[1]{*\eng{#1}}

\newcommand{\com}[1]{\hfill (\emp{#1})}%

\usepackage{relsize,xspace}
\newcommand{\into}{\ensuremath{\rightarrow}\xspace}
\newcommand{\tot}{\ensuremath{\leftrightarrow}\xspace}
\usepackage{url}
\newcommand{\lurl}[1]{\MyLogo{\url{#1}}}

\usepackage{langsci-gb4e}
\let\eachwordone=\slshape
%%% langsci-gb4e defines \thexnumi as \@roman\@xsi{xnumi}, which only works
%%% inside an exe environment; some decks use xlisti on its own.  Restore
%%% the plain gb4e definition so that keeps working.
\makeatletter
\def\thexnumi{\@xsi{xnumi}}
\makeatother
\newcommand{\lx}[1]{\textbf{\mtciteform{#1}}}
\newcommand{\ix}{\ex\slshape}



\newcommand{\ent}{\ensuremath{\Rightarrow}\xspace}
\newcommand{\ngv}{\ensuremath{\not\Rightarrow}\xspace}
%\usepackage{times}
%\usepackage{nttfoilhead}
\newcommand{\myslide}[1]{\foilhead[-25mm]{\raisebox{12mm}[0mm]{\emp{#1}}}\MyLogo{\logo}}
\newcommand{\myslider}[1]{\rotatefoilhead[-25mm]{\raisebox{12mm}[0mm]{\emp{#1}}}}
%\newcommand{\myslider}[1]{\rotatefoilhead{\raisebox{-8mm}{\emp{#1}}}}

\newcommand{\section}[1]{\myslide{}{\begin{center}\Huge \emp{#1}\end{center}}}



\usepackage[lyons,j,e,k]{mtg2e}
%\renewcommand{\mtcitestyle}[1]{\textcolor{-red!75!green!50}{\textsl{#1}}}
\renewcommand{\mtcitestyle}[1]{\textcolor{teal}{\textsl{#1}}}
\newcommand{\cs}{\mtciteform}
\newcommand{\iz}[1]{\texttt{\textup{#1}}}
\newcommand{\gm}{\textsc}
\usepackage[normalem]{ulem}
\newcommand{\ul}{\uline}
\newcommand{\ull}{\uuline}
\newcommand{\wl}{\uwave}
\newcommand{\vs}{\ensuremath{\Leftrightarrow}~}
%%%
%%% Bibliography
%%%
\usepackage{natbib}
%\usepackage{url}
\usepackage{bibentry}



\usepackage[utf8]{inputenc}
%%% trees
\usepackage{multicol}
%%% pst-node: one example in lec-12-control draws a connector between
%%% two points with \rnode/\ncbar.  This used to come in with rtrees.
\usepackage{pst-node}
\usepackage{forest}
\useforestlibrary{linguistics}
\forestapplylibrarydefaults{linguistics}

%%% \makebox's optional arguments are read as AVM delimiters by
%%% langsci-avm, so hide them behind a macro with no visible brackets.
\newcommand{\lbox}[1]{\makebox[0mm][l]{#1}}

%%% Tree styles.
%%%   words  -- trees whose leaves are actual words (italicised, aligned)
%%%   cat    -- trees whose leaves are categories or feature structures
\forestset{
  words/.style={for tree={s sep=7mm, l sep=5mm,
                          if n children=0{font=\itshape, tier=terminal}{}}},
  cat/.style={for tree={s sep=7mm, l sep=5mm}},
}
\newcommand{\wtree}[1]{\begin{forest} words #1 \end{forest}}
\newcommand{\ctree}[1]{\begin{forest} cat #1 \end{forest}}
%%% the rule being applied at this step of a derivation
\newcommand{\step}[1]{\begin{center}\emp{#1}\end{center}}

%%% AVMs (Language Science Press)
\usepackage{langsci-avm}

%%% Macros the decks use, carried over from headx.tex
\newcommand{\blank}{\rule{3em}{1pt}\xspace}
\newcommand{\ft}[1]{\textsc{#1}}
\newcommand{\idn}{\mtciteform}    %%% Indonesian
\newcommand{\prd}[1]{\textbf{#1}}
\newcommand{\typ}[1]{\textit{#1}}
%%% \val: provide before renewing, in case a package defines it.
\providecommand{\val}[1]{}
\renewcommand{\val}[1]{\textit{#1}}
%%% Manning's \avmbox, still used in tree labels and in maths where a
%%% bare langsci \tag would not work.  \avm{\tag{...}} gives the same
%%% boxed tag as the ones inside feature structures.
\newcommand{\avmbox}[1]{\mbox{\avm{\tag{#1}}}}

%%% Ring a value in a dotted oval.  These used pstricks \ovalnode,
%%% which came in with rtrees; tikz (via forest) does the same job.
\newcommand{\OV}[1]{\tikz[baseline]{\node[draw,dotted,red,ellipse,
          inner sep=1pt,outer sep=0pt]{#1};}}
\newcommand{\OVB}[1]{\tikz[baseline]{\node[draw,dotted,blue,ellipse,
          inner sep=1pt,outer sep=0pt]{#1};}}

%%% From CSLI book
\newcommand{\mc}{\multicolumn}
\newcommand{\HD}{\textbf{H}\xspace}
\newcommand{\el}{< >}       %%% empty list; langsci-avm needs the space

\newcommand{\A}{\noindent\textbf{A}: }
\newcommand{\Q}{\noindent\textbf{Q}: }

%%% Question and answer on two pages, written once:
%%%   \qaslide{Title}{... question ... \answer{... answer ...}}
%%% prints the slide twice, first without the \answer parts and then
%%% with them, so the answer appears "on click".  The example counter
%%% is reset so both pages show the same example numbers.
\newif\ifshowanswer
\newcommand{\answer}[1]{\ifshowanswer#1\fi}
\newcounter{qaxnumi}
\newcommand{\qaslide}[2]{%
  \setcounter{qaxnumi}{\value{xnumi}}%
  \showanswerfalse\myslide{#1}#2%
  \setcounter{xnumi}{\value{qaxnumi}}%
  \showanswertrue\myslide{#1}#2%
  \showanswerfalse}
%\newcommand{\C}{\noindent\textbf{C}: }



\begin{document}

\header{Lecture 5}{Semantics}{}
\maketitle


\myslide{Overview}
\begin{itemize}
\item The linguist's stance
\item Which aspects of meaning we model
\item Semantic features and principles: building the semantics of phrases
\item Modification and coordination
\item Structural ambiguity
\item Building a tree (updated)
\item Problems: two kinds of modifiers, number names
\end{itemize}

\section{Review}

\myslide{Last Week}
%\MyLogo{Chapter 4}

\begin{itemize}
\item Complex feature structures allow us to
  \begin{itemize}
  \item Write more general rules
  \item Constrain them lexically
  \end{itemize}
\end{itemize}

\myslide{Head-Complement Rule}

\avm{[\type*{phrase}  val & [ comps & <  > ] ]
              \ $\rightarrow$\ \HD \ 
             [\type*{word}
             val & [ comps & < \1, \ldots, \tag{n} > ] ]\ \ \1, \ldots, \tag{n}
}
\begin{itemize}
\item The possible complements are specified lexically \\
\begin{small}\begin{tabular}{@{}ll}
    \avm{< \textnormal{devour}, \ [\type*{word}
      head & \type{verb} \\
      % [\type*{agr-cat} agr & [ per & 3 \\ num & sg ] ]
      val & [ COMPS & < NP > ]
      ]>
} &
    \avm{< \textnormal{put}, \ [\type*{word}
      head & \type{verb} \\
      % [\type*{agr-cat} agr & [ per & 3 \\ num & sg ] ]
      val & [ COMPS & < NP PP > ]
      ]>
} 
  \end{tabular}
\end{small}
\end{itemize}

\myslide{Head-Specifier Rule}

\avm{[\type*{phrase}\\
             VAL & [ COMPS & < >\\
                      SPR & < >]]\ \ \
$\rightarrow$\ \ \ \ {\2}\ \ \ \ \HD[ %\type{phrase}\\
			                VAL & [ COMPS & < >\\
                                                 SPR & < {\2} > ]]
}
\begin{itemize}
\item Combines the rules expanding S and NP (and potentially other categories).
\item Again, restrictions on specifiers come from the lexicon.
\end{itemize}
\begin{small}\begin{tabular}{ll}
    \avm{< \textnormal{dog}, \ [\type*{word}
      head & \type{noun} \\
      % [\type*{agr-cat} agr & [ per & 3 \\ num & sg ] ]
      val & [ spr & < D > \\ comps & < > ]
      ]>
} &
    \avm{< \textnormal{eat}, \ [\type*{word}
      head & \type{verb} \\
      % [\type*{agr-cat} agr & [ per & 3 \\ num & sg ] ]
      val & [ spr & < NP > \\  comps & < (NP) > ]
      ]>
} 
  \end{tabular}
\end{small}

\myslide{Two Principles}
\begin{itemize}
\item \txx{The Valence Principle}
  \begin{quote}
    Unless the rule says otherwise, the mother’s
    values for the VAL features (SPR and
    COMPS) are identical to those of the head
    daughter.
  \end{quote}
\item \txx{The Specifier-Head Agreement Constraint (SHAC)}
\\ Verbs and common nouns must be specified as:
\begin{center}
  \avm{[ head & [ agr & \1 ] \\
    val & [ spr & < [ head & [ agr & \1 ] ] > ] ]
}
\end{center}
In plain English: a verb or noun agrees with its specifier
\end{itemize}

\myslide{Overview}
\begin{itemize}
\item The linguist's stance
\item Which aspects of meaning we model
\item Semantic features and principles: building the semantics of phrases
\item Modification and coordination
\item Structural ambiguity
\item Building a tree (updated)
\item Problems: two kinds of modifiers, number names
\end{itemize}


\section{The Linguist's Stance}

\myslide{The Linguist's Stance: Building a Precise Model}

\begin{itemize}
\item Some of our statements are statements about how the model
works:
\begin{quote}
  “[\val{prep}] and [\ft{agr} \val{3sing}] can’t be combined because \ft{agr} is not a
  feature of the type \val{prep}.”
\end{quote}
\item Some of our statements are statements about how (we think)
English or language in general works.
\begin{quote}
  “The determiners \eng{a} and \eng{many} only occur with count nouns, the
  determiner \eng{much} only occurs with mass nouns, and the determiner \eng{the}
  occurs with either.”
\end{quote}
\item Some are statements about how we code a particular
linguistic fact within the model.
\begin{quote}
  “All count nouns are [\ft{spr}  $<$ [\ft{count} +]$>$].”
\end{quote}
\end{itemize}

\myslide{The Linguist's Stance: A Vista on Possible Sentences}

We look at the set of possible English sentences
\begin{itemize}
\item \ldots  as a background against which linguistic
  elements (words, phrases) have a distribution
\item \ldots  as an arena in which linguistic elements
  “behave” in certain ways
\end{itemize}

\emp{In plain English: we learn what a word is by looking at where it
  can and can't appear.}


\section{Our Slice of Meaning}

\qaslide{Semantics: Where's the Beef?}{

So far, our grammar has no semantic representations. We
have, however, been relying on semantic intuitions in our
argumentation, and discussing semantic contrasts where
they line up (or don't) with syntactic ones.

\medskip
Examples?
\answer{
\begin{itemize}
\item \emp{structural ambiguity}
\item \emp{S/NP parallelism}
\item \emp{count/mass distinction}
\item \emp{complements vs.\ modifiers}
\end{itemize}}
}

\myslide{Our Slice of a World of Meanings}

Aspects of meaning we \emp{won't} account for (in this course)
\begin{itemize}
\item Pragmatics
\item Fine-grained lexical semantics
  \\ The meaning of \eng{life} is
  \begin{itemize}\addtolength{\itemsep}{2ex}
  \item \iz{life}  (or \textsf{life}$'$)
\item
  \avm{%\avmvalfont{\it}
    [ reln & life \\ inst & i ]
}
\item Not like WordNet: 
\iz{life$_1$}  $\subset$ \iz{being$_1$} $\subset$ \iz{state$_1$} \ldots

 % -- (a characteristic state or mode of living; "social life"; "city life"; "real life")
 %       => being, beingness, existence -- (the state or fact of existing; "a point of view gradually coming into being"; "laws in existence for centuries")
 %           => state -- (the way something is with respect to its main attributes; "the current state of knowledge"; "his state of health"; "in a weak financial state")
 %               => attribute -- (an abstraction belonging to or characteristic of an entity)
 %                   => abstraction, abstract entity -- (a general concept formed by extracting common features from specific examples)
 %                       => entity -- (that which is perceived or known or inferred to have its own distinct existence (living or nonliving))
 \end{itemize}

 \end{itemize}
\emp{We use the word itself as the name of its meaning: we don't try
  to define what \eng{life} is.}


% RELN life
% INST i

\myslide{Our Slice of a World of Meanings}
\noindent\avm{[mode &  prop \\
 index &  $s$\\
 restr & < [ RELN & {\bf save}\\
               SIT & $s$\\
              SAVER &  $i$\\
              SAVED &  $j$ ] \ , 
              [RELN & {\bf name}\\
              NAME &  Chris\\
              NAMED & \ $i$]\ , 
              [RELN & {\bf name}\\
              NAME &  Pat\\
              NAMED & \ $j$] > ]
} 
\begin{quote}
  “\ldots the linguistic meaning of \eng{Chris saved Pat} is a proposition
  that will be true just in case there is an actual situation that
  involves the saving of someone named Pat by someone named Chris.”
  \hfill \citep[p.~140]{sag:wasow:bender:2003}
\end{quote}
\emp{Read it as: there is a situation $s$ that is a saving; the saver
  $i$ is named Chris; the saved $j$ is named Pat.}

\myslide{Our Slice of a World of Meanings}

What we are accounting for is the \txx{compositionality} of
sentence meaning.

\begin{itemize}
\item How the pieces fit together
  \begin{quote}
    \txx{Semantic arguments} and \txx{indices}    
  \end{quote}
\item How the meanings of the parts add up to the meaning
of the whole.
\begin{quote}
  Appending \ft{restr} lists up the tree  
\end{quote}
\end{itemize}

The value of \txx{\ft{restr}iction} is the set of conditions that must
hold (in some possible world) for the expression to be applicable.

\myslide{Semantics in Constraint-Based Grammar}
\begin{itemize}
\item Constraints as (generalized) truth conditions
  \begin{itemize}
\item \txx{proposition}: what must be the case for a proposition to be true
  \hfill (\eng{Kim slept.})
\item \txx{directive}: what must happen for a directive to be fulfilled
  \hfill (\eng{Sleep!})
\item \txx{question}: the kind of situation the asker is asking about
  \hfill (\eng{Did Kim sleep?})
\item \txx{reference}: the kind of entity the speaker is referring to
  \hfill (\eng{the cat})
\end{itemize}
\item \txx{Syntax/semantics interface}: 
  \begin{quote}
    Constraints on how syntactic arguments are related to semantic
    ones, and on how semantic information is compiled from different
    parts of the sentence.
  \end{quote}
\end{itemize}

\section{Semantic Features and Principles}

\myslide{Feature Geometry}
\avm{[\type*{expression}
syn & [\type*{syn-cat}
      head & [\type*{pos} ... ] \\
      val & [ spr & < ... > \\
               comps & < ... > ]]\\
      sem & [\type*{sem-cat}
      mode &  \{ \val{prop}, \val{ques}, \val{dir}, \val{ref}, \val{none} \} \\
      index &  \{ $i, j, k, \ldots, s_1, s_2, \ldots$ \} \\
      restr & < ... > ]]
}



\myslide{SEM in Plain English}
\begin{tabular}{@{}lp{0.72\textwidth}}
  \ft{mode} & What kind of meaning: \val{prop} (a statement),
              \val{ques} (a question), \val{dir} (a command),
              \val{ref} (something referred to), or \val{none} \\[2ex]
  \ft{index} & A label for the thing or situation the expression is
               about: $i, j, k, \ldots$ for individuals, $s$ for
               situations \\[2ex]
  \ft{restr} & A list of facts (\txx{predications}) that must be true
               of it \\
\end{tabular}

\bigskip
Each predication has a \ft{reln} (what kind of fact) and
\txx{roles} that say who is involved, e.g.\
\avm{[ reln & {\bf sleep} \\ sit & $s$ \\ sleeper & $i$ ]}: in
situation $s$, $i$ sleeps

%%% Frame a value: used to pick out the parts of a tree under discussion.
\newcommand{\nd}[2]{{\setlength{\fboxsep}{1pt}\fcolorbox{#1}{white}{#2}}}

\myslide{How the Pieces Fit Together: A Name}
\avm{< \eng{Kim}, [syn & [head & [\type*{noun}\\
  agr & \type{3sing}]\\
  val & [spr & {< >}\\
  comps & {< >}]]\\
  sem & [mode & ref\\
  index & \ \ \nd{red}{$i$}\\
  restr & < [reln & {\bf name}\\
  name & Kim\\
  named & \ \ \nd{red}{$i$}] > ]] >
}

\begin{itemize}
\item \ft{index} $i$: the person this word is about
\item \ft{named} $i$: that same person is called Kim
\end{itemize}

\myslide{How the Pieces Fit Together: A Verb}
\avm{< \eng{slept} ,\ [syn & [head & \type{verb}\\
                                     val & [spr & < NP$_{\nd{red}{$i$}}$ >\\
                                     comps & <  >]]\\
                            sem & [mode & prop\\
                                    index & \nd{blue}{$s$}\\
                                    restr & < [reln & {\bf sleep}\\
                                                 sit & \ \ \nd{blue}{$s$}\\
                                                 sleeper & \ \ \nd{red}{$i$}  ] >]] >
}

\begin{itemize}
\item \ft{index} $s$: the sleeping situation
\item NP$_i$ is short for an NP whose \ft{index} is $i$
\item So whoever is the specifier (the subject) is the \ft{sleeper}
\end{itemize}

\myslide{The Pieces Together}

\begin{forest}
cat
[{S} [{\avm{\1} NP  \avm{[ sem & [ index & $i$ ] ]}} [{\eng{Kim}}]]
[{VP  \avm{[ syn & [ val & [ spr & < \1 > ] ] \\
      sem & [ mode & prop\\
      index & $s$\\
      restr & < [reln & {\bf sleep}\\
      sit & \ \ $s$\\
      sleeper & \ \ $i$  ] > ] ]}} [{\eng{slept}}]]]
\end{forest}

\begin{itemize}
\item The Head-Specifier Rule makes the NP the VP's specifier \avm{\1}
\item So \lex{Kim}'s index $i$ is the \ft{sleeper}
\end{itemize}

\myslide{How to Share Semantic Information}

We need two more principles
\begin{itemize}
\item \txx{The Semantic Inheritance Principle (SIP)}
  \begin{quote}
    In any headed phrase, the mother's \ft{mode} and
    \ft{index} are identical to those of the head daughter.
  \end{quote}
  \emp{A phrase is about what its head is about.}
\item \txx{The Semantic Compositionality Principle (SCP)}
  \begin{quote}
    In any well-formed phrase structure, the mother's
    \ft{restr} value is the sum of the \ft{restr} values of
    the daughters.
  \end{quote}
  \emp{A phrase means everything its parts mean.}
\end{itemize}

\myslide{List Sum}

The SCP adds lists together with $\oplus$ (technically, concatenation):
just join the lists, in order

\begin{center}
  \avm{< a, b > } $\oplus$ \avm{< c > } $=$ \avm{< a, b, c >}
\end{center}

Order matters:
\begin{center}
  \avm{< c > } $\oplus$ \avm{< a, b > } $=$ \avm{< c, a, b >}
\end{center}

In general:
$\langle a_1, \ldots, a_n \rangle \oplus \langle b_1, \ldots, b_m \rangle =
\langle a_1, \ldots, a_n,  b_1, \ldots, b_m \rangle $

%%% Kim slept, with the S node filled in step by step.  #1, #2 and #3
%%% are the S node's MODE, INDEX and RESTR values; #4 is the text below.
\newcommand{\kimtree}[4]{%
\begin{center}\scalebox{0.75}{\begin{forest}
cat
[{S  \avm{[ sem & [ mode & #1 \\ index & #2 \\ restr & #3 ] ]}}
  [{\avm{\1} NP  \avm{[ sem & [ mode & ref \\ index & $i$\\
      restr & < [reln & {\bf name}\\
      name & \ \ Kim \\ named & \ \ $i$ ] >  ] ]}} [{\eng{Kim}}]]
  [{\HD VP  \avm{[ syn & [ val & [ spr & < \1 > ] ] \\
      sem & [  mode & prop \\ index & $s$ \\
      restr & < [reln & {\bf sleep}\\
      sit & \ \ $s$\\
      sleeper & \ \ $i$  ] > ] ]}} [{\eng{slept}}]]]
\end{forest}}\end{center}
#4}

\myslide{Filling in the S Node}
\kimtree{\blank}{\blank}{\blank}{%
  \lex{Kim} and \lex{slept} get their \ft{sem} from their lexical
  entries.  S has no lexical entry: where do its \ft{mode}, \ft{index}
  and \ft{restr} come from?}

\myslide{Filling in the S Node: SIP}
\kimtree{\nd{red}{prop}}{\nd{red}{$s$}}{\blank}{%
  \emp{SIP}: VP is the head daughter, so S has the same \ft{mode} and
  \ft{index} as the VP}

%%% An AVM can't be nested inside another one's value, so typeset the
%%% framed RESTR list first.
\newsavebox{\scpbox}
\sbox{\scpbox}{\nd{red}{\avm{< [reln & {\bf name}\\
      name & \ \ Kim  \\ named & \ \ $i$ ] ,
      [reln & {\bf sleep}\\
      sit & \ \ $s$\\
      sleeper & \ \ $i$  ]  >}}}

\myslide{Filling in the S Node: SCP}
\kimtree{prop}{$s$}{\usebox{\scpbox}}{%
  \emp{SCP}: S's \ft{restr} is the NP's \ft{restr} $\oplus$ the VP's
  \ft{restr}}

%%% The cat slept on the mat, following the index i.  #1--#4 are the
%%% frame colours of the links in the chain (white = invisible), so
%%% every page has the same layout.  #5 is the text on the right.
\newcommand{\cattree}[5]{%
\noindent\begin{minipage}[t]{0.64\textwidth}\vspace{0pt}
\scalebox{0.7}{\begin{forest}
cat
[{S} [{\avm{\1} \nd{#3}{NP$_i$}} [{D} [{\eng{the}}]]
  [{\nd{#4}{Nom$_i$}} [{\eng{cat}}]]]
[{VP \nd{#2}{\avm{[ spr & < \1 > ]}}}
  [{VP \avm{[ spr & < \nd{#1}{\avm{\1}NP$_i$} > \\
          restr & < [reln & {\bf sleep}\\
          sit & \ \ $s$\\
          sleeper & \ \ \nd{#1}{$i$}  ] >  ]}} [{\eng{slept}}]]
  [{PP} [{\eng{on the mat}}, roof]]]]
\end{forest}}
\end{minipage}\hfill
\begin{minipage}[t]{0.34\textwidth}\vspace{0pt}\small\raggedright
\textbf{Why is \lex{cat}'s index the \ft{sleeper}?}
#5
\end{minipage}}

\myslide{What Identifies Indices?}
\cattree{white}{white}{white}{white}{}

\myslide{What Identifies Indices?}
\cattree{red}{white}{white}{white}{%
  \begin{enumerate}
  \item \emp{\lex{slept}'s lexical entry}: its specifier's index
    is the \ft{sleeper}
  \end{enumerate}}

\myslide{What Identifies Indices?}
\cattree{white}{red}{white}{white}{%
  \begin{enumerate}
  \item \lex{slept}'s lexical entry: its specifier's index
    is the \ft{sleeper}
  \item \emp{Valence Principle}: the Head-Modifier Rule doesn't
    change \ft{spr}, so the upper VP has the same \ft{spr}
    \avm{< \1 >}
  \end{enumerate}}

\myslide{What Identifies Indices?}
\cattree{white}{white}{red}{white}{%
  \begin{enumerate}
  \item \lex{slept}'s lexical entry: its specifier's index
    is the \ft{sleeper}
  \item Valence Principle: the upper VP has the same \ft{spr}
    \avm{< \1 >}
  \item \emp{Head-Specifier Rule}: the NP is that specifier
    \avm{\1}, so its index is $i$
  \end{enumerate}}

\myslide{What Identifies Indices?}
\cattree{white}{white}{white}{red}{%
  \begin{enumerate}
  \item \lex{slept}'s lexical entry: its specifier's index
    is the \ft{sleeper}
  \item Valence Principle: the upper VP has the same \ft{spr}
    \avm{< \1 >}
  \item Head-Specifier Rule: the NP is that specifier
    \avm{\1}, so its index is $i$
  \item \emp{SIP}: the NP gets its \ft{index} from its head, the
    Nom \lex{cat}
  \end{enumerate}}

\myslide{Answer: A Chain of Identities}
\cattree{red}{red}{red}{red}{%

  \medskip
  A chain of identities makes the \ft{index} of \lex{cat} the
  \ft{sleeper} in the \textbf{sleep} predication of \lex{slept}.

  \medskip
  \emp{The cat is the one who slept.}}

\myslide{Summary: Words \ldots }
\begin{itemize}
\item contribute predications \hfill (\textbf{sleep})
\item ‘expose’ one index in those predications, for use by words or phrases
  \hfill (\emp{\ft{index} $s$ = \ft{sit} $s$})
\item relate syntactic arguments to semantic arguments
  \hfill (\blu{NP$_i$ = \ft{sleeper} $i$})
 \\ \scalebox{0.8}{\avm{< \eng{slept} ,\ [syn & [head & \type{verb}\\
                                     val & [spr & < NP$_{\nd{blue}{$i$}}$ >\\
                                     comps & <  >]]\\
                            sem & [mode & prop\\
                                    index & \nd{red}{$s$}\\
                                    restr & < [reln & {\bf sleep}\\
                                                 sit & \ \ \nd{red}{$s$}\\
                                                 sleeper & \ \ \nd{blue}{$i$}  ] >]] >
}}
\end{itemize}

\myslide{Summary: Grammar Rules \ldots }
\begin{itemize}
\item Identify feature structures (including the \ft{index} value) across daughters
  \begin{itemize}\addtolength{\itemsep}{1ex}
  \item Head-Specifier Rule \\
\scalebox{0.9}{\avm{[\type*{phrase} 
  val & [ spr & < > ] ] \ \ \
  $\rightarrow$\ \ \ \ {\1}\ \ \ \ \HD[ %\type{phrase}\\
  val & [ spr & < {\1} >  \\ comps & < > ]]
}}
\item Head-Complement Rule \\
\scalebox{0.9}{\avm{[\type*{phrase}
    val & [ comps & <  > ] ]
    \ $\rightarrow$\ \HD \ 
    [\type*{word}
    val & [ comps & < \1, \ldots, \tag{n} > ] ]\ \ \1, \ldots, \tag{n}
}}
\item Head-Modifier Rule \\
\scalebox{0.9}{\avm{[\type*{phrase} ]
    \ $\rightarrow$\ \HD \ \1  [ val & [ comps & < > ]  ] \ \ 
    [ val & [ comps & < > \\ mod & < \1 > ] ]
}}
\end{itemize}
\end{itemize}

\myslide{Summary: Grammar Rules \ldots }

\begin{itemize}
\item Identify feature structures (including the \ft{index} value) across daughters
\item License trees that are subject to the semantic principles
  \begin{itemize}
  \item SIP: ‘passes up’ \ft{mode} and \ft{index} from head daughter
  \item SCP: ‘gathers up’ predications (\ft{restr} list) from all daughters
  \end{itemize}
\item The semantics is strictly compositional --- all of the meaning
  comes from the words, rules and principles.
  \begin{itemize}
  \item We then enrich this with pragmatic inference --- but we need a
    base to infer from
  \end{itemize}
\end{itemize}

\section{Modification and Coordination}

\myslide{Other Aspects of Semantics}
\begin{itemize}
\item Tense, quantification (only touched on here)
\item Modification
\item Coordination
\item Structural ambiguity
\end{itemize}

\myslide{Evolution of a Phrase Structure Rule}
\begin{itemize}
\item [C2] NOM \into NOM PP; VP \into VP PP
\item [C3] 
\scalebox{0.8}{\avm{[\type*{phrase} val & [ spr & - \\ comps & itr ] ] 
    \ $\rightarrow$\ \HD \  [\type*{phrase}  val & [ spr & - ]   ] \ \ 
    PP
}}
\item [C4] 
\scalebox{0.8}{\avm{[\type*{phrase}  ]
    \ $\rightarrow$\ \HD \ [  val & [ comps & < > ]  ] \ \ 
    PP
}}
\item[C5]
\scalebox{0.8}{\avm{[\type*{phrase} ]
    \ $\rightarrow$
    \HD \ \1  [ syn &  [ val & [ comps & < > ]  ] ] \ \ 
    [ syn & [ val & [ comps & < > \\ mod & < \1 > ] ] ]
}}

\item[=] 
\scalebox{0.8}{\avm{[\type*{phrase} ]
     \ $\rightarrow$\ \HD \ \1   [ comps & < >  ] \ \ 
      [ comps & < > \\ mod & < \1 > ]
}} 
\end{itemize}

\myslide{Evolution of Another Phrase Structure Rule}

\begin{itemize}
\item [C2] X \into X+ CONJ X
\item [C3] 
\scalebox{0.8}{\avm{\1 \ \into \ \1+ \ \ 
    [\type*{word} head & conj ] \ \ 
    \1
}}
\item [C4] 
\scalebox{0.8}{\avm{[ val & \1 ] \ \into \ [ val & \1 ]+ \ \ 
    [\type*{word} head & conj ] \ \ 
    [ val & \1 ]
}}
\item[C5]
\scalebox{0.8}{\avm{[ syn & [ val & \0 ] \\ 
       sem & [ ind & $s_0$ ]  ]
    \ \into \
}} \\
\scalebox{0.65}{\avm{[ syn & [ val & \0 ] \\ 
       sem & [ ind & $s_1$ ]  ] \ \ldots \  
   [ syn & [ val & \0 ] \\ 
     sem & [ ind & $s_{n-1}$ ] ]
    [ syn & [ head & conj ] \\
    sem & [ ind  & $s_0$ \\ 
    restr & < [ args & < $s_1, \ldots, s_{n-1}, s_n$ > ] > ]  ]\ \ 
  [ syn & [ val & \0 ] \\ 
     sem & [ ind & $s_n$ ] ]
}}
\item[=]
\scalebox{0.75}{\avm{[  val & \0 \\ ind & $s_0$ ]
    \ \into \ 
  [  val & \0 \\ ind & $s_1$  ] \ \ldots \  
   [  val & \0  \\ ind & $s_{n-1}$ ]
    [ head & conj \\
     ind  & $s_0$ \\ 
    restr & < [ args & < $s_1, \ldots, s_{n-1}, s_n$ > ] >  ]\ \ 
  [  val & \0 \\ ind & $s_n$ ]
}}
\end{itemize}
% \item[=] 
% \scalebox{0.8}{\begin{avm}\avmfont{\sc}
%      \[ \asort{phrase} \]
%      \ $\rightarrow$\ \HD \ \@1   \[ comps & < >  \] \ \ 
%       \[ comps & < > \\ mod & \< \@1 \> \] 
%  \end{avm}} 

\myslide{Combining Constraints and Coordination}

\begin{itemize}
\item  Coordination Rule \\
\scalebox{0.75}{\avm{[  val & \0 \\ ind & $s_0$ ]
    \ \into \ 
  [  val & \0 \\ ind & $s_1$  ] \ \ldots \  
   [  val & \0  \\ ind & $s_{n-1}$ ]
    [ head & conj \\
     ind  & $s_0$ \\ 
    restr & < [ args & < $s_1, \ldots, s_{n-1}, s_n$ > ] >  ]\ \ 
  [  val & \0 \\ ind & $s_n$ ]
}}
\item Lexical Entry for \lex{and} \\
  \scalebox{0.9}{\avm{< \eng{and} ,\ [ syn & [head & \type{conj} ] \\
                          sem & [mode & none\\
                                  ind & $s$\\
                                  restr & < [reln & {\bf and}\\
                                               sit & \ \ $s$\\ ] >]] >
}}
\end{itemize}

\myslide{Combining Constraints and Coordination}
\noindent\scalebox{0.68}{\begin{forest}
cat
[{S \avm{[ ind & $s_0$ \\
restr & 
< [reln & {\bf name} \\ name & Joe \\ named & $j$ ],
 [reln & {\bf joke} \\ sit & $s_1 $ \\ joker & $j$ ], 
\0 \ [reln & {\bf and}\\
sit & \ \ $s_0$\\ 
args & < $s_1, s_2$ > ], 
[reln & {\bf name} \\ name & Kim \\ named & $k$], 
[reln & {\bf smile} \\ sit & $s_2 $ \\ smiler & $k$] > ]}} [{S \avm{[ ind & $s_1$ ]}} [{\eng{Joe jokes}}, roof]]
[{\avm{[ head & \type{conj} \\
         ind & $s_0$\\
        restr & < \0 [reln & {\bf and} ] > ]}} [{\eng{and}}]]
[{S \avm{[ ind & $s_2$ ]}} [{\eng{Kim smiles}}, roof]]]
\end{forest}}

\myslide{How to Read the Coordination Rule}
\begin{itemize}
\item The rule is \emp{not headed}, so the SIP doesn't apply: the rule
  itself has to say what the mother's \ft{ind} is
  \begin{itemize}
  \item The SCP still applies: the mother's \ft{restr} is all the
    daughters' \ft{restr}s added together
  \end{itemize}
\item \avm{\0}: the mother and all the conjuncts have the same
  \ft{val} --- you can only coordinate things that need the same
  things
\item The mother's \ft{ind} $s_0$ is the conjunction's \ft{ind}
\item The conjunction's predication links the conjuncts: its \ft{args}
  are their indices $\langle s_1, \ldots, s_n \rangle$
\end{itemize}

\qaslide{Where Does Each Index Come From?}{
In \eng{Joe jokes and Kim smiles}:
\begin{itemize}
\item $s_1$ and $s_2$? \answer{\emp{From \lex{jokes} and \lex{smiles},
    passed up to each S by the SIP}}
\item $s_0$? \answer{\emp{From the lexical entry for \lex{and}; the
    Coordination Rule makes it the mother's \ft{ind}}}
\item \ft{args} $\langle s_1, s_2 \rangle$? \answer{\emp{The
    Coordination Rule: the indices of the two conjuncts}}
\item The whole \ft{restr} list? \answer{\emp{The SCP: everything the
    daughters' \ft{restr}s contain, in order}}
\end{itemize}
}

\section{Structural Ambiguity}

\myslide{Structural Ambiguity}
\eng{Joe jokes and Kim smiles often} has two readings:
\begin{itemize}
\item[I] \lex{often} modifies the whole coordination \\
  \emp{It is often the case that Joe jokes and Kim smiles}
\item[II] \lex{often} modifies only \eng{Kim smiles} \\
  \emp{Joe jokes, and Kim often smiles}
\end{itemize}

\bigskip
Our grammar gives each reading its own tree, and the two trees
build different \ft{restr} lists

%%% The RESTR list of the top S of both trees; #1 is often's ARG.
\newcommand{\jokerestr}[1]{%
\noindent\scalebox{0.6}{\avm{[ ind & $s_0$ \\
mode & prop \\
restr &
< [reln & {\bf name} \\ name & Joe \\ named & $j$ ],
 [reln & {\bf joke} \\ sit & $s_1 $ \\ joker & $j$ ],
 [reln & {\bf and}\\
sit & \ \ $s_0$\\
args & < $s_1, s_2$ > ],
[reln & {\bf name} \\ name & Kim \\ named & $k$],
[reln & {\bf smile} \\ sit & $s_2 $ \\ smiler & $k$],
[reln & {\bf often} \\ sit & $s_3 $ \\ arg & \nd{red}{#1} ]  > ]}}}

\myslide{Tree I: \textit{often} Modifies the Coordination}
\begin{center}
\scalebox{0.85}{\begin{forest}
cat
[{S \avm{[ ind & $s_0$ ]}}
  [{\avm{\1} S \avm{[ ind & \nd{red}{$s_0$} ]}}
    [{S \avm{[ ind & $s_1$ ]}} [{\eng{Joe jokes}}, roof]]
    [{\avm{[ head & \type{conj} \\ ind & $s_0$ ]}} [{\eng{and}}]]
    [{S \avm{[ ind & $s_2$ ]}} [{\eng{Kim smiles}}, roof]]]
  [{ADV \avm{[ mod & < \1 > ]}} [{\eng{often}}, roof]]]
\end{forest}}
\end{center}
The top S:\\
\jokerestr{$s_0$}

\myslide{Tree II: \textit{often} Modifies \textit{Kim smiles}}
\begin{center}
\scalebox{0.85}{\begin{forest}
cat
[{S \avm{[ ind & $s_0$ ]}}
  [{S \avm{[ ind & $s_1$ ]}} [{\eng{Joe jokes}}, roof]]
  [{\avm{[ head & \type{conj} \\ ind & $s_0$ ]}} [{\eng{and}}]]
  [{S \avm{[ ind & $s_2$ ]}}
    [{\avm{\1} S \avm{[ ind & \nd{red}{$s_2$} ]}} [{\eng{Kim smiles}}, roof]]
    [{ADV \avm{[ mod & < \1 > ]}} [{\eng{often}}, roof]]]]
\end{forest}}
\end{center}
The top S:\\
\jokerestr{$s_2$}

\qaslide{Question About Structural Ambiguity}{
Why isn’t this a possible semantic representation for
the string \eng{Joe jokes and Kim smiles often}?

\medskip
\jokerestr{$s_1$}

\answer{\par\medskip
\emp{\lex{often}'s \ft{arg} is $s_1$, the index of \eng{Joe jokes}.
  But \lex{often} can only modify a phrase it combines with, and the
  only ones next to it are \eng{Kim smiles} ($s_2$) and the whole
  coordination ($s_0$). No tree puts \lex{often} with \eng{Joe jokes}
  alone.}}
}

\myslide{Some Standard Extensions}

\begin{itemize}
\item Quantification
  \begin{itemize}
  \item typically expressed as restrictions on scope
  \item Minimal Recursion Semantics goes further
  \end{itemize}
\item Pragmatics
  \begin{itemize}
  \item typically expressed as another feature: \ft{context}
  \item contains things like \val{speaker, hearer, audience}
  \item used for pronominal reference, politeness
\end{itemize}
\end{itemize}

\myslide{What About Scope?}
\eng{Every dog chases a cat} has two readings:
\begin{itemize}
\item Each dog chases a cat, maybe a different cat for each dog
\item There is one particular cat, and every dog chases it
\end{itemize}
The difference is which quantifier is ``on top'' (has \txx{wide scope})

\begin{itemize}
\item We won't deal with scope in this course
\item Implemented HPSG grammars (with Minimal Recursion Semantics)
  leave it \txx{underspecified}: one representation that covers both
  readings
\item Usually we don't need to decide which reading is meant
\end{itemize}


\section{Building a Tree (Updated)}

%%% A type name with its features underneath.  The features are a
%%% bracketed column: langsci-avm mismeasures a column of bare features.
\newcommand{\tn}[2]{\begin{tabular}{@{}c@{}}#1\\
  $\left[\mbox{\scshape\begin{tabular}{@{}l@{}}#2\end{tabular}}\right]$\end{tabular}}

\myslide{Chapter 5 Type Hierarchy}
\MyLogo{New in Chapter 5 in red}
\begin{center}
  \scalebox{0.68}{\begin{forest}
cat, for tree={font=\itshape, s sep=3mm}
[{feat-struc} [{\tn{expression}{\emp{syn, sem}}} [{word}]
[{phrase}]]
[{\tn{\emp{syn-cat}}{\emp{head} \\ \emp{val}}}]
[{\tn{\emp{sem-cat}}{\emp{mode} \\ \emp{index} \\ \emp{restr}}}]
[{\emp{predication}}]
[{\tn{agr-cat}{per, num}} [{non-3sing} [{non-1sing} [{2sing}]
[{pl}]]
[{1sing}]]
[{\tn{3sing}{gend}}]]
[{\tn{val-cat}{spr \\ comps \\ \emp{mod}}}]
[{pos} [{\tn{agr-pos}{agr}} [{\tn{noun}{case}}]
[{\tn{verb}{aux}}]
[{\tn{det}{count}}]]
[{prep}]
[{adj}]
[{\emp{adv}}]
[{conj}]]]
\end{forest}}
\end{center}

\myslide{Algorithm for Building a Tree (Updated)}
\small
\textbf{Create the words (leaf nodes)}
\begin{enumerate}
\item Copy all the information from each word's lexical entry,
  \emp{including its \ft{sem}}
\item Enforce the SHAC (on verbs and common nouns)
\item Check the type hierarchy and add every appropriate feature
  (leave the value blank if it is underspecified)
\end{enumerate}
\textbf{Create the phrases}
\begin{enumerate}\setcounter{enumi}{3}
\item Find a rule that fits; add everything it says
\item Apply the principles (HFP, Valence Principle, \emp{SIP, SCP}) to
  update features
\item Repeat 4--5 until you reach the S node
\item Check: every feature you left blank should now have a value
\end{enumerate}
\normalsize

\section{Problems}

\myslide{Problem 5.1: Two Kinds of Modifiers in English}
In English, modifiers of nouns can appear either
before or after the noun, although any given modifier is usually
restricted to one position or the other.

\begin{itemize}
\item[(i)] \eng{The red dog on the roof}
\item[(ii)]  \bad{The on the roof dog}
\item[(iii)]  \bad{The dog red}
\end{itemize}
%
Our current Head-Modifier Rule only licenses post-head modifiers
(like \eng{on the roof} in (i)).

\myslide{Problem 5.1: Two Kinds of Modifiers (cont.)}
\begin{itemize}
\item[A.] Write a second Head-Modifier Rule that licenses
\index{pre head modifier@pre-head modifier}pre-head modifiers
(e.g., \eng{red} in (i)).
\item[B.] Modify the Head-Modifier 1 and Head-Modifier 2 Rules so that
they are sensitive to which kind of modifier is present and don't
generate (ii) or (iii).  [\textsl{Hint: Use a feature
{\rm [}POST-HEAD \{+,$-$\}{\rm ]} to distinguish \eng{red} and \eng{on the roof}.}]
\item[C.] Is POST-HEAD a \index{HEAD}HEAD feature? Why or why not?
\item[D.] Give lexical entries for \eng{red} and \eng{on} that show
the value of POST-HEAD. (You may omit the SEM features in these entries.)
\item[E.] Is (i) ambiguous according to your grammar (i.e.\
the Chapter 5 grammar modified to include the two Head-Modifier Rules,
instead of just one)? Explain your answer.
\end{itemize}

\myslide{Problem 5.1: Two Kinds of Modifiers (cont.)}
This problem assumed that we don't want to make the two Head-Modifier
Rules sensitive to the part of speech of the modifier.  One reason for
this is that modifiers of the same part of speech can occur before and
after the head, even though individual modifiers might be restricted
to one position or the other.

\begin{itemize}
\item[F.] Provide three examples of English NPs with
\index{adjective}adjectives or APs after the noun.
\item[G.] Provide three examples of \index{adverb (ADV)}adverbs that can come before
the verbs they modify.
\item[H.] Provide three examples of adverbs that can come after
the verbs they modify.
\index{ambiguity}
\end{itemize}


\myslide{Problem 5.3: Semantics of Number Names}
In  Chapter 3, we considered the
syntax of English number names, and in particular how
to find the head of a number name expression. Based on the
results of that problem, the lexical entry for
\eng{hundred} in a number name like \eng{two hundred five}
should include the constraints in (i). (Here we are assuming
a new subtype of \val{pos}, \val{number}, which is appropriate
for number name words.)
\index{pos@{\it pos}}
\index{number@{\it number}}

\begin{itemize}
\item[(i)] \scalebox{0.85}{\avm{< \textnormal{\eng{hundred}} , [ SYN [ HEAD & \type{number}\\
                       VAL & [ SPR & < [ HEAD & \type{number} ] >\\
                                COMPS & < [ HEAD & \type{number} ] > ]]] >
}}
\end{itemize}

\myslide{Problem 5.3: Semantics of Number Names (cont.)}
This lexical entry interacts with our ordinary Head-Complement
and Head-Specifier Rules to give us the phrase structure shown in
(ii):

\begin{itemize}
\item[(ii)] \scalebox{0.8}{\begin{forest}
cat, for tree={l sep=3mm}
[{NumP} [{NumP} [{\eng{two}}]]
[{Num$'$} [{Num} [{\eng{hundred}}]]
[{NumP} [{\eng{five}}]]]]
\end{forest}}
\end{itemize}

\myslide{Problem 5.3: Semantics of Number Names (cont.)}
Smith (1999) provides a compositional semantics
of number names. The semantics of this NP should be (iii):

\begin{itemize}
\item[(iii)]
\scalebox{0.55}{\avm{[ INDEX & $i$\\
   MODE & ref\\
   RESTR & < [ RELN & {\bf constant}\\
		 INST & $l$\\
		 VALUE & 2],
              [ RELN & {\bf times}\\
		 RESULT & $k$\\
		 FACTOR1 & $l$\\
                 FACTOR2 & $m$ ],
	     [ RELN & {\bf constant}\\
		 INST & $m$\\
		 VALUE & 100 ],
              [ RELN & {\bf plus}\\
 		 RESULT & $i$\\
	         TERM1 & $j$\\
		 TERM2 & $k$ ],
              [ RELN & {\bf constant}\\
		 INST & $j$\\
                 VALUE & 5 ] > ]
}}
\end{itemize}

\noindent
This  expresses ``(two times one hundred) plus five'' (i.e.\
205) as a FS.

\myslide{Problem 5.3: Semantics of Number Names (cont.)}
\begin{itemize}
\item[A.] Assume that the two constant predications
with the values 2 and
5 are contributed by the lexical entries for \eng{two} and \eng{five}.
What predications must be on the \index{RESTR}RESTR list of the lexical entry for
\eng{hundred} in order to build (iii) as the SEM value of \eng{two
hundred five}?
\item[B.] The lexical entry for \eng{hundred} will identify the
indices of its specifier and complement with the value of some feature
of a predication on its RESTR list.  Which feature of which predication
is the \index{index}index of the specifier identified with?  What about the index
of the complement?
\item[C.] The lexical entry for \eng{hundred} will identify its
own INDEX with the value of some feature of some predication on
its RESTR list.  Which feature of which predication must this
be, in order for the grammar
to build (iii) as the SEM value of \eng{two hundred five}?
\end{itemize}

\myslide{Problem 5.3: Semantics of Number Names (cont.)}
\begin{itemize}
\item[D.] Based on your answers in parts (A)--(C),
give a lexical entry for \eng{hundred} that includes
the constraints in (i) and a fully specified SEM value.
[{\sl Note: Your lexical entry need only account for \eng{hundred}
as it is used in \eng{two hundred five}. Don't worry about other
valence possibilities, such as
\eng{two hundred}, \eng{two hundred and five}, or \eng{a hundred}.}]
\item[E.] The syntax and semantics of number names do
not line up neatly: In the syntax, \eng{hundred} forms a constituent
with \eng{five}, and \eng{two} combines with \eng{hundred five} to
give a larger constituent.  In the semantics, the constant
predications
with the values 2 and 100 are related via the times
predication.
The result of that is related to the constant
predication
with the value
5, via the plus
predication.
Why is this mismatch not a problem for the grammar?
\end{itemize}

\myslide{Summary}
\begin{itemize}
\item We model how the meanings of the parts add up
  (\txx{compositionality}), not pragmatics or fine word meanings
\item \ft{sem}: \ft{mode} (what kind of meaning), \ft{index} (what it
  is about) and \ft{restr} (the facts that must hold)
\item Words link their syntactic arguments to semantic roles
  (NP$_i$ = \ft{sleeper} $i$)
\item SIP: a phrase is about what its head is about \\
  SCP: a phrase means everything its parts mean
\item Coordination has no head, so its rule states the semantics
\item Structurally ambiguous strings get different \ft{restr} lists
\end{itemize}

\myslide{Acknowledgments and References}

\begin{itemize}
\item Course design and slides borrow heavily from Emily Bender's course:
\textit{Linguistics 566: Introduction to Syntax for Computational Linguistics}
\\ \url{http://courses.washington.edu/ling566}
\item This lecture complements \citet[Chapter 5]{sag:wasow:bender:2003}
\end{itemize}

%\input{lq-04-semantics}

\bibliographystyle{aclnat}
\bibliography{abb,mtg,nlp,ling}
\end{document}


%%% Local Variables: 
%%% coding: utf-8
%%% mode: latex
%%% TeX-PDF-mode: t
%%% TeX-engine: xetex
%%% End: 

